The amplitude is the largest displacement from the equilibrium (the line $y = 0$), the period the time from one crest to the next. With $y(t) = A\cos(\omega t - \varphi_0)$ the crests come where $\omega t - \varphi_0 = 0, 2\pi, \dots$: the first crest $t_1$ after $t = 0$ gives $\varphi_0 = \omega\,t_1$, shifted by $2\pi$ into the range $-\pi < \varphi_0 \le \pi$. As a check: $y(0) = A\cos\varphi_0$, and $y$ rises at $t = 0$ if $\varphi_0 > 0$ ($\sin\varphi_0>0$), falls if $\varphi_0 < 0$.

\begin{abcliste}
\abc
$\result{A = \AaO}$, $\result{T = \TaO}$, so
\begin{align}
 \omega &= \waF = \frac{2\pi}{\Ta} = \wa
\end{align}
The first crest is at $t_1 = \tcaO$, a quarter period:
\begin{align}
 \varphi_0 &= \phaF = \wa\cdot\tca = \pha = \result{\frac{\pi}{2}}
\end{align}
Check: at $t = 0$ the graph is at $y = 0$ and rising.
\begin{align}
 y(t) &= \AaO\cdot\cos\left(\frac{2\pi}{\TaO}\,t - \frac{\pi}{2}\right)
\end{align}
(the same as $\AaO\cdot\sin(\omega t)$).

\abc
$\result{A = \AbO}$, $\result{T = \TbO}$, so
\begin{align}
 \omega &= \wbF = \frac{2\pi}{\Tb} = \wb
\end{align}
The first crest is at $t_1 = \tcbO$, five eighths of a period:
\begin{align}
 \omega\,t_1 &= \wb\cdot\tcb = \phbr = \frac{5\pi}{4}
\end{align}
This is larger than $\pi$, so subtract $2\pi$:
\begin{align}
 \varphi_0 &= \frac{5\pi}{4}-2\pi = \phb = \result{-\frac{3\pi}{4}}
\end{align}
Check: $y(0) = A\cos\varphi_0 = \ybP{-2}{2}$, and the graph falls at $t = 0$, as it must for $\varphi_0 < 0$.
\begin{align}
 y(t) &= \AbO\cdot\cos\left(\frac{2\pi}{\TbO}\,t + \frac{3\pi}{4}\right)
\end{align}
\end{abcliste}
